\input{e.tex}

\corrPolitique{0}

\newcommand{\eps}{\varepsilon}

\usepackage{url}

\title{TP 2 de Calcul Scientifique 2\\ Méthode des Eléments Finis}

\author{Master MIGS 2}

\date{2020/2021}

\begin{document}

\maketitle

%Le but de ce TP est d'apprendre à utiliser un code éléments finis "dernière génération", qui se base sur l'écriture du problème à résoudre sous forme faible.
%
%On utilisera dans ce but FreeFEM++, qui est un logiciel libre développé principalement par le Laboratoire Jacques Louis Lions, à Paris 6 (Frédéric Hecht est le principal contributeur).
%On pourra compléter cette utilisation, si besoin, avec des logiciels de visualisation.

%La première partie, sur l'équation de Poisson, a principalement pour objectif d'illustrer le cours et les notions les plus importantes, et aussi de permettre une prise en main du logiciel. La deuxième partie,
Ce deuxième TP porte sur l'équation de la chaleur, et a pour objectif d'appliquer les notions vues en cours à un problème "réel".

\parindent0mm
\section{Equation de la chaleur}

On va résoudre numériquement l'équation de la chaleur, d'abord sur un exemple simple, mais artificiel, et ensuite pour un problème "réel" (thermique d'une aile d'avion).

\subsection{Un premier exemple simple}

Soit $T>0$ et \[
\Omega := (0,1) \times (0,1).
\]
On va d'abord résoudre l'équation de la chaleur sur le domaine $\Omega \times [0,T]$ de $\eR^3$:
\[
\partial_t u -\Delta u = f \textrm{ dans } \Omega \times [0,T],
\quad
u = 0 \textrm{ sur } \partial \Omega \times [0,T],
\quad
u(t=0) = 0 \textrm{ dans } \Omega,
\]
avec $f \in L^2(\Omega \times (0,T))$ le terme source, donné, et où l'inconnue est $u : \Omega \times [0,T] \rightarrow \eR$. On prendra, en fait, ici $f=1$ (on chauffe uniformément sur tout le carré).

\begin{enumerate}

\item Rappeler quelle est l'expression de la formulation totalement discrète en utilisant à la fois des éléments finis $\mathbb{P}_1$ pour la variable d'espace et le schéma d'Euler implicite pour la variable temporelle. 

\item Ecrire le script FreeFEM++ correspondant. Visualiser la solution obtenue (choisir, par exemple, comme pas de temps et d'espace $h=0.1$, $\tau = 0.1$).

\item A partir de quelle valeur de $T$ observez-vous une solution (quasi-)stationnaire ?

\item Etudier comment converge la méthode lorsque $h$ et $\tau$ deviennent petit. Quelle est la meilleure stratégie pour faire évoluer $h$ et $\tau$ ?

%\item Qu'est-ce que vous observez si vous prenez le schéma d'Euler explicite pour la discrétisation en temps ?

\end{enumerate}

{\it Remarque.} outre les scripts fournis avec FreeFEM++, vous pouvez aussi reprendre le script fourni sur la page web de Grégoire Allaire : 

\begin{center}
\url{http://www.cmap.polytechnique.fr/~allaire/map431/infoX2.html}
\end{center}

{\it Questions subsidiaires:

\begin{enumerate}

\item Que se passe-t'il si on arrête de chauffer au bout d'un temps donné ? Pouvez-vous fournir une explication mathématique ?

\item Arrivez-vous à gagner en précision avec d'autres schémas en temps et d'autres éléments finis en espace ?

\end{enumerate}

}

\subsection{Application en aéronautique}

Simulons maintenant l'évolution de la température au sein d'une aile d'avion en acier, représentée Figure~\ref{fig:airfoil}. 
On considère que l'aile est à une température initiale 
$u_0=298.15^{\circ}{\rm K}$, et qu'elle est entourée d'air, qui, lui, est à une température $u_\mathrm{s}=253.15^{\circ}{\rm K}$ (sur la partie $\Gamma_{\mathrm{R}}$). 
On impose une température $u_{{d}} = 348.15^{\circ}{\rm K}$ sur une petite portion 
$\Gamma_{\mathrm{D}}$ de la frontière, en bas de l'aile ($\partial \Omega = 
\Gamma_{\mathrm{R}} \cup \Gamma_{\mathrm{D}})$. %, see Figure~\ref{fig:airfoil} (a)).
L'acier est caractérisé par sa densité volumique $\rho=7860\  \mbox{kg/m}^3$, 
une chaleur spécifique $c=502 \ \mbox{J/Kg\,K}$ et une conductivité thermique $\kappa=60\mbox{W/m\,K}$. 
Pour l'air, le coefficient de transfert de chaleur par convection est $\alpha=45\ \mbox{W/m$^2$\,K}$. On prendra un temps de simulation $T=1000$ s. La dimension transverse $L$ de l'aile est $0.8$ m.
%
\begin{figure}[ht!]
\center
\includegraphics[width=0.6\linewidth]{airfoil3.pdf} %(a)
\caption{Géométrie de l'aile d'avion, avec notations pour le domaine et les bords.}
\label{fig:airfoil}
\end{figure}
%
Le modèle mathématique pour ce problème est le suivant :
\begin{equation*}
%\label{eq:heatAirfoil}
\begin{aligned}
\rho c\ \partial_t {u} - \kappa\:\Delta u &= 0 &\quad\textrm{dans}&\quad\Omega\times(0,T),\\
u(\cdot,0)&=u_0 &\quad\textrm{dans}&\quad\Omega,\\
\kappa \frac{\partial u}{\partial n}+\alpha u&=\alpha u_s & \quad\textrm{sur}&\quad \Gamma_{\mathrm{R}},\\
u & = u_{{d}} & \quad\textrm{sur}&\quad \Gamma_{\mathrm{D}}.
\end{aligned}
\end{equation*}
%
Ecrire le script FreeFEM++ qui permet de résoudre le problème décrit ci-dessus, puis prédire et visualiser l'évolution de la température au niveau de l'aile.

\end{document}

\newpage

\section{Technique de pénalité}

Autre sujet, à voir en fonction des avancées.

Reprendre matériel précédent (exams CTU et mémoires).

\end{document}

